\section{A continuous section for the Frobenius comparison}
\label{app:frobenius}

The solid Frobenius comparison used in
\Cref{prop:constant-acyclicity} requires surjectivity on continuous
families, not only on point values.  We supply the topological step for
its precise function ring.  This avoids applying a linear open-mapping
theorem to Frobenius minus identity, which is not linear over the
algebraically closed coefficient field.

Let $C$ be an algebraically closed complete nonarchimedean field of
characteristic $p$, containing $k=\F_q$, and let
\[
 R=H^0(\widetilde{\mathbb B}^{\,*}_C,\cO)
   =\left\{\sum_{\alpha\in\Z[1/p]}c_\alpha t^\alpha:
                 \text{the series converges for every }0<|t|<1\right\}.
\]
Its Fr\'echet topology is given by the Gauss seminorms
$\|f\|_r=\sup_\alpha |c_\alpha|r^\alpha$, $0<r<1$.
Convergence means that these weighted coefficients form a $c_0$ family
for every such radius.  Here $\phi_q$ acts on the $C$-coefficients
and fixes $t$; it is not the absolute Frobenius of $R$.

\begin{lemma}\label{lem:continuous-as-section}
The additive map $\phi_q-1:R\to R$ has a continuous section as a map
of sets.  Its condensification is surjective, and its kernel, with the
induced topology, is the valued-field completion
$E=\bigl(\bigcup_{j\ge0}k((t^{1/p^j}))\bigr)^{\wedge}$.
Consequently the canonical map $E\to R^{h\phi_q}$ is an equivalence
of derived solid $\F_p$-modules.
\end{lemma}

\begin{proof}
For each coset of the open additive subgroup $C^{\circ\circ}$ in $C$,
choose a representative $a$ and a root $b_a$ of $b_a^q-b_a=a$.
Choose $a=b_a=0$ on the zero coset.  Define
\[
 s(a+\epsilon)=b_a-\sum_{j\ge0}\epsilon^{q^j},
 \qquad |\epsilon|<1.
\]
The sum converges, and $s(x)^q-s(x)=x$.
If $|x-y|<1$, the points lie in the same coset and the first term
of the difference series dominates, so $|s(x)-s(y)|=|x-y|$.
If $|x-y|\ge1$, the equation for $s(x)-s(y)$ gives norm
$|x-y|^{1/q}$ (including norm one at the boundary).
Thus $s$ is continuous, $s(0)=0$, and
\[
 |s(x)-s(y)|\le\max\{|x-y|,|x-y|^{1/q}\}.
\]
Define $S$ coefficientwise.  For every $0<r<1$ this estimate gives
\[
 \|S(f)-S(g)\|_r
 \le\max\{\|f-g\|_r,\ \|f-g\|_{r^q}^{1/q}\}.
\]
It also gives the corresponding estimate on every coefficient tail.
Since $r^q$ is again in $(0,1)$, $S$ preserves convergence and is
continuous for the Fr\'echet topology.  It is a section of $\phi_q-1$.
For any profinite space $T$, postcomposition with $S$ lifts every
continuous map $T\to R$.  This proves the asserted surjectivity of
condensed groups, and therefore of solid modules.

The kernel consists of the same convergent series with coefficients in
$k$.  Its nonzero coefficients all have norm one.  The convergence
condition says that the exponents in its support form a left-finite
subset of $\Z[1/p]$: below any bound there are only finitely many.
Truncation therefore identifies it with the completed perfection $E$.
On this kernel every Gauss seminorm is $r^{v_t(f)}$; hence its induced
topology is exactly the valued-field topology of $E$.
The two-term complex $[R\xrightarrow{\phi_q-1}R]$ now has kernel
$E$ and zero cokernel in solid modules.  It computes homotopy fixed
points for the infinite cyclic action, proving the last assertion.
\end{proof}

For the application, the quasi-Stein perfectoid calculation in
\cite[Lemma~3.6.2]{BMR2026} identifies the relevant solid cohomology
with this function ring.  The lemma proves the enhancement needed in
\cite[Corollary~3.6.8]{BMR2026} directly for that ring.
The equivalence is the canonical map from constants; its construction
as a map is independent of the chosen section.  It therefore commutes
with all actions and base changes already commuting with that map.
No equivariance of the auxiliary set-theoretic section is required.
